\documentclass[12pt,a4paper]{article}

% ===== PACKAGE IMPORTS =====
\usepackage[utf8]{inputenc}
\usepackage[english]{babel}
\usepackage[a4paper,left=17mm,top=22mm,right=17mm,bottom=15mm]{geometry}

% Math packages
\usepackage{amsmath,amsfonts,amssymb}
\usepackage{mathtools}
\usepackage{bm}
\usepackage[T1]{fontenc}
\usepackage{amsthm}

% Graphics and floats
\usepackage{graphicx}
\usepackage{float}
\usepackage{booktabs}

% Bibliography and citations
\usepackage{natbib}

% Links and references
\usepackage[colorlinks=true, citecolor=blue, linkcolor=black, urlcolor=blue]{hyperref}
\usepackage{url}

% Code listings
\usepackage{listings}
\usepackage{xcolor}

% Configure Python syntax highlighting
\lstdefinelanguage{Python}{
	keywords={and, as, assert, break, class, continue, def, del, elif, else, 
		except, False, finally, for, from, global, if, import, in, is, 
		lambda, None, nonlocal, not, or, pass, raise, return, True, try, 
		while, with, yield, async, await},
	keywordstyle=\color{blue}\bfseries,
	ndkeywords={self, cls, int, float, str, list, dict, set, tuple, bool, range},
	ndkeywordstyle=\color{purple}\bfseries,
	sensitive=true,
	comment=[l]{\#},
	commentstyle=\color{gray}\itshape,
	string=[b]",
	stringstyle=\color{red},
	basicstyle=\ttfamily\footnotesize,
	breaklines=true, breakatwhitespace=true,
	tabsize=4, showstringspaces=false
}

\lstset{
	language=Python,
	frame=single,
	numbers=left,
	numberstyle=\tiny\color{gray},
	xleftmargin=2em, framexleftmargin=1.5em,
	basicstyle=\ttfamily\small,
	commentstyle=\color{gray},
	keywordstyle=\color{blue},
	stringstyle=\color{red},
	backgroundcolor=\color{gray!10},
	rulecolor=\color{black!30},
	captionpos=b,
	breaklines=true, breakatwhitespace=false,
	tabsize=4, showstringspaces=false
}


\providecommand{\keywords}[1]{\textbf{\textit{Keywords---}} #1}

% ===== THEOREM STYLES =====
\theoremstyle{plain}
\newtheorem{theorem}{Theorem}[section]
\newtheorem{proposition}[theorem]{Proposition}
\newtheorem{lemma}[theorem]{Lemma}
\newtheorem{corollary}[theorem]{Corollary}

\theoremstyle{definition}
\newtheorem{definition}[theorem]{Definition}
\newtheorem{assumption}[theorem]{Assumption}

\theoremstyle{remark}
\newtheorem{remark}[theorem]{Remark}

% ===== SMALL MACROS =====
\newcommand{\R}{\mathbb{R}}
\newcommand{\E}{\mathbb{E}}
\newcommand{\Vol}{\mathrm{Vol}}
\newcommand{\Tr}{\mathrm{Tr}}
\newcommand{\diag}{\mathrm{diag}}
\newcommand{\supp}{\mathrm{supp}}
\newcommand{\1}{\mathbf{1}}

\title{%
	\textbf{Epiplexity And Graph Wiring\\
		An Empirical Study\\ for the design of a generic algorithm}
}
\author{Lorenzo Moriondo\\
	\small Independent Researcher --- tuned.org.uk\\
	\small ORCID: 0000-0002-8804-2963}

\date{15 March 2026}

% ════════════════════════════════════════════════════════════════════
\begin{document}
	\maketitle
	
	% ─── Abstract ───────────────────────────────────────────────────────
	\begin{abstract}
		Having introduced \emph{Graph Wiring}---a technique that leverages the
		Graph Laplacian computed in feature space to provide semantically-aware
		search over high-dimensional vector corpora---and \emph{MRR-Top0} as a
		topology-aware retrieval metric for evaluating it; this paper proceeds
		to demonstrate formally and empirically that the feature-space Laplacian
		produced by ArrowSpace carries \emph{structural information} in the sense
		of epiplexity \citep{finzi2026epiplexity}, so that it is not plain graph
		metadata but a semantic artifact.
		Using the CVE~1999--2025 vulnerability corpus as a case study, we
		instantiate ArrowSpace as a spectral vector-search engine, wrap its
		feature-space Laplacian in a Laplacian-constrained Gaussian Markov
		Random Field (LGMRF), and evaluate the resulting two-part Minimum
		Description Length (MDL) code.
		The model achieves a compression ratio of~$38.4\times$ over raw float32
		storage, passes all three structural-information diagnostic tests. Furthermore
		\emph{it is demonstrated that the same Laplacian object as computed by Graph Wiring supports six distinct algorithm families
		(search, classification, anomaly detection, diffusion, dimensionality
		reduction, and data valuation) without additional learning}.
		The two accompanying Jupyter notebooks are intended as a reproducible
		reference pattern for applying epiplexity measurement in algorithm design
		for large-scale data engineering for LLMs and ML operations.
	\end{abstract}
	
	\keywords{Epiplexity, Graph Laplacian, Graph Wiring, ArrowSpace,
		Spectral indexing, MDL, Structural information, Vector search,
		CVE, LGMRF}
	
	% ===================================================================
	\section{Introduction}
	\label{sec:intro}
	
	\paragraph{Graph Wiring}
	\cite{graphwiring2025} introduced Graph Wiring as a principled way to
	turn any matrix of high-dimensional embeddings into a discrete graph
	whose nodes are \emph{features} (dimensions), not items.
	Edges connect features that co-vary systematically across the dataset,
	and the resulting \emph{feature-space Laplacian}~$L_F$ governs a
	Dirichlet energy that measures how spectrally coherent each item vector
	is on the feature manifold.
	The ArrowSpace library~\citep{Moriondo2025ArrowSpace} implements Graph
	Wiring as a production-ready spectral index, exposing a
	\emph{taumode}~$\lambda$-score per item that blends Rayleigh-quotient
	energy with cosine-based geometry.
	
	\paragraph{MRR-Top0}
	\cite{mrrtop02025} extended Mean Reciprocal Rank to the full top-$k$
	list via Personalised PageRank, conductance, and modularity, creating a
	retrieval metric sensitive to the topological structure of ranked lists.
	A key observation in that paper is that MRR-Top0 admits a natural
	information-theoretic interpretation: the difference between
	MRR-Top0 and standard MRR measures the structural information captured
	by spectral re-ranking that a position-unaware metric ignores.
	This connection motivates the present study.
	
	\paragraph{This paper}
	We ask three sharply-stated questions:
	\begin{enumerate}
		\item Does the ArrowSpace feature-space Laplacian carry
		non-trivial \emph{structural information} in the sense of
		epiplexity~\citep{finzi2026epiplexity}?
		\item How much information does the ArrowSpace computation itself
		\emph{generate}, relative to raw embedding storage?
		\item What are the practical implications of this analysis for data
		engineering in LLM and machine-learning pipelines?
	\end{enumerate}
	To answer these questions empirically we produce two Jupyter notebooks:\\
	(i)~\texttt{00\_arrowspace\_epiplexity\_structural\_information.ipynb},
	a self-contained tutorial establishing the epiplexity machinery on
	synthetic data; and\\ (ii)~\texttt{01\_arrowspace\_cve1999\_2025\-%
		epiplexity\_check\_v3.ipynb}, a full case study on the CVE~1999--2025
	corpus using a production ArrowSpace index.
	
	Full code used is available \href{https://github.com/tuned-org-uk/graph-wiring-epiplexity}{in this repository}
	
	The remainder of the paper is structured as follows.\\
	Section~\ref{sec:background} reviews epiplexity and ArrowSpace.\\
	Section~\ref{sec:prelim} walks through the tutorial notebook.\\
	Section~\ref{sec:notebook} presents the CVE case-study notebook.\\
	Section~\ref{sec:results} discusses results and SOTA comparisons.\\
	Section~\ref{sec:conclusion} concludes.
	
	% ===================================================================
	\section{Background}
	\label{sec:background}
	
	\subsection{Epiplexity and time-bounded MDL}
	
	Classical information theory conflates randomness and structure.
	Shannon entropy $H(X)$ increases identically whether one adds
	independent noise or discovers new regularities.
	Epiplexity, introduced by \cite{finzi2026epiplexity}, resolves this
	by splitting information into \emph{structural} and \emph{random} parts
	for a computationally bounded observer.
	
	\begin{definition}[\cite{finzi2026epiplexity}, Def.\ 7]
		A prefix-free program~$P$ is a \emph{$T$-time probabilistic model} over
		$\{0,1\}^n$ if it supports:
		\begin{itemize}
			\item \textbf{Evaluation:} $\mathrm{Prob}_P(x) = P(0, x)$ computes
			in time $T(n)$;
			\item \textbf{Sampling:} $\mathrm{Sample}_P(u) = P(1, u)$ draws
			$x\sim P$ in time $T(n)$;
			\item \textbf{Normalisation:} $\sum_{x}\mathrm{Prob}_P(x)=1$.
		\end{itemize}
	\end{definition}
	
	\begin{definition}[\cite{finzi2026epiplexity}, Def.\ 8]
		The \emph{epiplexity} $S_T(X)$ and \emph{time-bounded entropy}
		$H_T(X)$ are defined via the minimising model:
		\[
		P^* = \arg\min_{P\in\mathcal{P}_T}
		\Bigl(|P| + \E\bigl[\log_2\tfrac{1}{P(X)}\bigr]\Bigr),
		\]
		\[
		S_T(X) = |P^*|, \qquad
		H_T(X) = \E\bigl[\log_2\tfrac{1}{P^*(X)}\bigr].
		\]
		$S_T(X)$ is the \emph{structural} information and $H_T(X)$ the
		\emph{random} information for a $T$-bounded observer.
	\end{definition}
	
	Operationally, one uses the two-part MDL code \citep{GrunwaldMDL,
		VitanyiMDL}:
	\begin{equation}
		\label{eq:mdl}
		\mathrm{MDL}_T(X)
		= \underbrace{|P^*|}_{\text{model bits}=S_T}
		+ \underbrace{\sum_{i=1}^{N} H_T(x_i)}_{\text{data bits}}.
	\end{equation}
	The \emph{compression test} for structural content is:
	\[
	\mathrm{MDL}_T(X) < N\cdot F\cdot b + O(1)
	\;\Longrightarrow\;
	P^*\text{ captures structural information},
	\]
	where $N\cdot F\cdot b$ is the uncompressed baseline at $b$~bits per
	coordinate \citep{VitanyiStructureFunction}.
	A key distinction from classical compression is that both sides of the
	inequality are \emph{observer-relative}: raising the time bound~$T$
	(or equivalently, the compute budget of the algorithm) reveals more
	structure, increasing $S_T$ and decreasing $H_T$.
	
	\subsection{Graph Wiring and the ArrowSpace Algorithm}
	
	Graph Wiring \citep{graphwiring2025} constructs an $F\times F$ Laplacian
	on the feature-space of an $N\times F$ embedding matrix by:
	
	\begin{enumerate}
		\item Clustering items into $C_0$ centroids and projecting to an
		$F$-dimensional centroid matrix $\mathbf{C}^{C_0\times F}$;
		\item Transposing $\mathbf{C}$ so that each \emph{column}
		(feature dimension) becomes a node, and building a $k$-NN
		graph with Gaussian edge weights to form $L_F = D - W$;
		\item Computing per-item Rayleigh quotients
		$R(x_i) = x_i^\top L_F x_i / x_i^\top x_i$;
		\item Mapping $R(x_i)$ through a taumode transform to a score
		$\lambda_i\in[0,1]$ used as a spectral index.
	\end{enumerate}
	
	The signal-on-graph interpretation is key \citep{ShumanDSPGraphs,
		SandryhailaDSPGraphs}: each item vector $x_i\in\R^F$ is a signal
	defined on the feature-node graph.
	A signal that is \emph{smooth} on $L_F$ (low Rayleigh quotient)
	corresponds to an item whose feature profile is consistent with the
	dominant correlations in the corpus --- in other words, a
	semantically typical item.
	A spectrally rough item is either anomalous or from a domain not
	well-represented in the index.
	ArrowSpace \citep{Moriondo2025ArrowSpace} implements this pipeline in
	Rust, exposing Python bindings and supporting Parquet-native vector
	stores.
	
	% ===================================================================
	\section{Epiplexity Computation: A Tutorial Notebook}
	\label{sec:prelim}
	
	Before presenting the CVE case study, this section walks through\\
	\texttt{00\_arrowspace\_epiplexity\_structural\_information.ipynb}, a
	self-contained notebook that introduces epiplexity measurement from
	scratch using synthetic data and a pure-Python reimplementation of the
	key ArrowSpace constructs.
	The notebook is structured in nine sections, mirroring the theoretical
	framework of~\cite{finzi2026epiplexity} and the pipeline stages of
	ArrowSpace \citep{Moriondo2025ArrowSpace}.
	
	\subsection{Cell 1 --- Dependencies}
	
	The tutorial requires only standard scientific Python libraries:
	
	\begin{lstlisting}[caption={Tutorial notebook: dependency cell},
	label={lst:deps00}]
import math, warnings
import numpy as np
import scipy.sparse as sp
import scipy.sparse.linalg as spla
from scipy.spatial.distance import cdist
from sklearn.neighbors import kneighbors_graph
import plotly.graph_objects as go
import plotly.express as px

warnings.filterwarnings('ignore')
np.random.seed(42)
print("All dependencies loaded successfully.")
	\end{lstlisting}
	
	Deliberately, no ArrowSpace import appears here; the tutorial implements
	its own feature-space Laplacian builder so that every mathematical step
	is transparent.
	
	\subsection{Cell 2 --- Core Epiplexity Definitions (§1)}
	
	The markdown cell recapitulates Definitions~7 and~8 from
	\cite{finzi2026epiplexity} and derives the two-part MDL
	code~\eqref{eq:mdl}.
	The following code cell then instantiates the formula numerically on
	three hypothetical dataset archetypes:
	
	\begin{lstlisting}[caption={Two-part MDL toy illustration},
		label={lst:mdltoy}]
def two_part_mdl(model_bits, per_item_entropies):
	return model_bits + sum(per_item_entropies)
	
datasets = {
	"Random noise (API keys, hashes)":
	{"model_bits": 8,   "mean_entropy": 35.0, "n": 1000},
	"Structured text (code, articles)":
	{"model_bits": 950, "mean_entropy": 12.0, "n": 1000},
	"LLM embeddings (ArrowSpace)":
	{"model_bits": 420, "mean_entropy":  9.5, "n": 1000},
}

for name, d in datasets.items():
	mdl = two_part_mdl(d['model_bits'],
	[d['mean_entropy']] * d['n'])
	raw = d['n'] * 32
	ratio = mdl / raw
	print(f"{name:<38} MDL={mdl:.0f} bits  ratio={ratio:.3f}x")
	\end{lstlisting}
	
	\noindent\textbf{Output:}
	\begin{lstlisting}[numbers=none, backgroundcolor=\color{white},
		frame=none]
		Random noise (API keys, hashes)    MDL=35008  ratio=1.094x
		Structured text (code, articles)   MDL=12950  ratio=0.405x
		LLM embeddings (ArrowSpace)        MDL= 9920  ratio=0.310x
	\end{lstlisting}
	
	The central insight is visible immediately: random data (ratio $>1$)
	\emph{cannot} be compressed by any model; structured text and LLM
	embeddings both compress substantially, with LLM embeddings achieving
	the lowest ratio because their semantic geometry is highly regular.
	The model-bits term $|P^*|=S_T$ is small but non-zero, capturing the
	cost of encoding the regularity itself.
	
	\subsection{Cell 3 --- The ArrowSpace Pipeline as a Prefix-Free Program (§2)}
	
	The §2 markdown cell introduces ArrowSpace's four-stage pipeline as a
	formal prefix-free program $P_\mathrm{AS}$ and derives its description
	length using Elias-gamma integer coding \citep{GrunwaldMDL}:
	
	\begin{equation}
		\label{eq:pas}
		|P_\mathrm{AS}|
		= \sum_{x\in\{N,F,C_0,k\}}\!\!(2\lfloor\log_2 x\rfloor+1)
		+ \underbrace{64+8+32}_{\text{seed, flags, }\sigma}
		+ \underbrace{C_0 F_\mathrm{eff} b}_{\text{centroids}}
		+ \underbrace{F k(\lceil\log_2 F\rceil + b)}_{\text{Laplacian topology}}.
	\end{equation}
	
	The code cell implements this exactly:
	
	\begin{lstlisting}[caption={Description length of the ArrowSpace program},
		label={lst:desclength}]
def elias_gamma_bits(x: int) -> int:
	return 2 * math.floor(math.log2(max(1, x))) + 1
		
def compute_description_length(N, F, C0=None, k=16, b=32, F_eff=None):
	if C0 is None:
		C0 = max(100, min(2000, int(2 * math.sqrt(N))))
	if F_eff is None:
		F_eff = F
	header   = sum(elias_gamma_bits(x) for x in [N, F, C0, k])
	params   = 64 + 8 + 32
	centroid = C0 * F_eff * b
	topology = F * k * (math.ceil(math.log2(max(2, F))) + b)
	total    = header + params + centroid + topology
	raw      = N * F * b
	return {"total_KB": total/8/1024, "raw_KB": raw/8/1024,
		"compression_ratio": raw/total, "C0": C0}
		
scales = [
	(1_000,     768,  "Small (1k / 768d)"),
	(100_000,  1536,  "Medium (100k / 1536d)"),
	(1_000_000,1536,  "LLM-scale (1M / 1536d)"),
	(1_000_000,4096,  "LLM-large (1M / 4096d)"),
]
for N, F, label in scales:
	r = compute_description_length(N, F)
	print(f"{label:<28} |P_AS|={r['total_KB']:.1f} KB "
	f"raw={r['raw_KB']:.0f} KB  {r['compression_ratio']:.1f}x")
	\end{lstlisting}
	
	\noindent\textbf{Output:}
	\begin{lstlisting}[numbers=none, backgroundcolor=\color{white},
		frame=none]
		Small (1k / 768d)            |P_AS|=363.0 KB  raw=3000 KB    8.3x
		Medium (100k / 1536d)        |P_AS|=3921.0 KB raw=600000 KB  153.0x
		LLM-scale (1M / 1536d)       |P_AS|=12129.0 KB raw=6000000 KB 494.7x
		LLM-large (1M / 4096d)       |P_AS|=32352.0 KB raw=16000000 KB 494.6x
	\end{lstlisting}
	
	This is a foundational result: the description length of the ArrowSpace
	\emph{model} is sublinear in $N$ because it depends on $C_0 \sim
	O(\sqrt{N})$ centroids, not on all~$N$ items directly.
	At LLM-scale ($10^6$ items, $1536$~dimensions), the structural program
	costs only $12$~MB while the raw embeddings consume~$6$~GB---a
	structural compression of nearly $500\times$ even before accounting for
	per-item entropy.
	
	\subsection{Cell 4 --- The LGMRF Probabilistic Wrapper (§3)}
	
	The §3 markdown cell identifies a conceptual gap: ArrowSpace computes
	deterministic $\lambda$-scores, but Definition~7 requires a model that
	can \emph{evaluate probabilities} and \emph{sample}.
	The fix is to wrap $L_F$ in a Laplacian-constrained Gaussian Markov
	Random Field (LGMRF) \citep{DongLaplacianGMRF, YingLaplacianConstraints,
		RueHeldGMRF}:
	
	\begin{equation}
		\label{eq:lgmrf}
		Q = \beta L_F + \gamma I \succ 0,\quad
		x\sim\mathcal{N}(0, Q^{-1}),\quad
		\log P_\mathrm{AS}(x) = -\tfrac{1}{2}x^\top Qx - \log Z.
	\end{equation}
	
	The \emph{algebraic bridge} between ArrowSpace's internal computation
	and the probabilistic model is:
	\[
	x^\top Qx
	= \beta\underbrace{x^\top L_F x}_{\text{Dirichlet energy}}
	+ \gamma\|x\|^2.
	\]
	The Dirichlet energy that ArrowSpace already computes is, up to a
	constant, the negative log-probability under the LGMRF.
	The following code class \texttt{ArrowSpaceProbabilisticModel}
	implements all three requirements of Definition~7:
	
	\begin{lstlisting}[caption={LGMRF wrapper satisfying Definition~7},
		label={lst:lgmrfclass}]
class ArrowSpaceProbabilisticModel:
	def __init__(self, L_F, beta=1.0, gamma=1e-3):
		self.F  = L_F.shape[0]
		self.L_F = L_F
		self.Q  = beta * L_F.tocsc() + gamma * sp.eye(self.F,
		format='csc')
		self._lu      = spla.splu(self.Q)
		self.log_det_Q = float(np.sum(
		np.log(np.abs(self._lu.U.diagonal()))))
		self.log_Z = (0.5 * self.F * np.log(2*np.pi)
		- 0.5 * self.log_det_Q)
		
	# Definition 7 - Evaluation  [O(nnz(L_F))]
	def evaluate_log_prob(self, x):
		return float(-0.5 * x @ (self.Q @ x) - self.log_Z)
		
	# Definition 7 - Sampling  [O(F^1.5)]
	def sample(self, n=1, rng=None):
		rng = rng or np.random.default_rng(0)
		z = rng.standard_normal((self.F, n))
		return self._lu.solve(z)
		
	# Definition 7 - Normalisation  [O(1) given log Z]
	def log_partition(self):
		return self.log_Z
		
	def time_bounded_entropy(self, x):
		return -self.evaluate_log_prob(x) / np.log(2)
		
	def rayleigh_quotient(self, x):
		d = float(x @ x)
		return float(x @ (self.L_F @ x)) / d if d > 1e-12 else 0.0
	\end{lstlisting}
	
	A quick sanity check on a 12-node path-graph Laplacian confirms the
	core theorem: smooth signals (constant vector, Dirichlet energy~$= 0$)
	obtain strictly higher probability than rough signals (alternating
	$\pm 1$, maximum Dirichlet energy):
	
	\begin{lstlisting}[numbers=none, backgroundcolor=\color{white},
		frame=none]
		Smooth  | Dirichlet energy: 0.0000 | log P: -11.979 | H_T: 17.28 bits
		Rough   | Dirichlet energy: 3.6667 | log P: -13.812 | H_T: 19.93 bits
		=> Smooth signal has HIGHER probability (lower H_T) -- LGMRF confirmed.
	\end{lstlisting}
	
	\noindent The two-bit gap between smooth and rough signals is the
	Laplacian's direct contribution to the epiplexity budget.
	
	\subsection{Cell 5 --- The Full MDL Toolkit and Compression Test (§4)}
	
	Having established the LGMRF model, the §4 cell assembles
	\texttt{ArrowSpaceMDLToolkit}, which combines the description-length
	encoder with the LGMRF to evaluate Equation~\eqref{eq:mdl} on any
	$N\times F$ dataset:
	
	\begin{lstlisting}[caption={MDL toolkit: compression test over three synthetic datasets},
		label={lst:mdltoolkit}]
class ArrowSpaceMDLToolkit:
	def __init__(self, X, k=5, beta=1.0, gamma=1e-3,
		sigma=1.0, n_centroids=None):
		self.N, self.F = X.shape
		self.C0  = n_centroids or max(10, min(200,
		int(2*np.sqrt(self.N))))
		self.L_F = build_knn_feature_laplacian(X, k=k, sigma=sigma)
		self.model = ArrowSpaceProbabilisticModel(
		self.L_F, beta=beta, gamma=gamma)
		self.rayleigh = np.array(
		[self.model.rayleigh_quotient(x) for x in X])
		self.entropy  = np.array(
		[self.model.time_bounded_entropy(x) for x in X])
		
	@property
	def structural_bits(self):
		return self.model.description_length_bits(self.C0, self.k)
		
	@property
	def compression_ratio(self):
		return (self.N * self.F * 32) / self.mdl_total
		
	def compression_test(self):
		passes = self.mdl_total < self.N * self.F * 32
		return {"compression_ratio": self.compression_ratio,
			"passes_compression": passes, ...}
	\end{lstlisting}
	
	Running the toolkit on three synthetic datasets---structured manifold
	(two Gaussian clusters), pure Gaussian noise, and a low-rank smooth
	manifold plus noise---yields the following results:
	
	\begin{table}[H]
		\centering
		\caption{MDL compression test on three synthetic datasets ($N=200$, $F=20$).}
		\label{tab:synthetic-mdl}
		\begin{tabular}{lrrrrl}
			\toprule
			Dataset & $|P_\mathrm{AS}|$ (KB) & $\sum H_T$ (KB)
			& MDL$_T$ (KB) & Raw (KB) & Pass? \\
			\midrule
			Structured manifold & 2.56 & 3.08 & 5.64 & 15.62 & \checkmark \\
			Pure random noise   & 2.56 & 3.08 & 5.64 & 15.62 & \checkmark \\
			Smooth low-rank     & 2.56 & 3.06 & 5.62 & 15.62 & \checkmark \\
			\bottomrule
		\end{tabular}
	\end{table}
	
	All three datasets pass the compression test, establishing that
	the \emph{model cost} $|P_\mathrm{AS}|$ alone---independent of
	data-specific entropy---accounts for the majority of the compression.
	This is a property of the quadratic normalisation constant $\log Z$
	in the LGMRF: a denser graph (higher $k$) produces a more informative
	model and a smaller $\log Z$, which in turn lowers $H_T$ per item.
	
	\subsection{Cell 6 --- Structural Diagnostics (§5)}
	
	Three independent tests formalise the
	\emph{metadata vs.\ structural information} decision:
	
	\begin{table}[H]
		\centering
		\caption{Three diagnostic tests for structural information content.}
		\label{tab:diagnostics}
		\begin{tabular}{lll}
			\toprule
			Test & Criterion & Interpretation \\
			\midrule
			Compression ratio  & $\mathrm{MDL}_T < n_\mathrm{raw}$
			& Model compresses data $\Rightarrow$ structural \\
			Spectral gap       & $\lambda_2(L_F) \gg 0$
			& Graph connected $\Rightarrow$ encodes real topology \\
			Rayleigh CV        & $\sigma(\lambda)/\mu(\lambda) > 0.5$
			& Items vary on manifold $\Rightarrow$ discriminative \\
			\bottomrule
		\end{tabular}
	\end{table}
	
	The \texttt{StructuralInformationDiagnostics} class runs all three tests
	and emits a \textsc{structural} or \textsc{metadata} verdict.
	On the synthetic datasets, the toy Laplacians (built on only $F=20$
	dimensions with a small-$N$ graph) pass compression but fail the
	spectral gap test---alerting the practitioner that a denser
	feature graph is needed.
	This is the intended pedagogical outcome of the tutorial: the tests
	jointly act as a sensitivity dashboard pointing to which hyperparameter
	needs adjustment.
	
	\subsection{Cells 7--9 --- Multi-Class Applications, LLM Tools, and Visualisations (§6--§9)}
	
	The remaining cells demonstrate that the same $L_F$, once computed,
	drives six independent downstream operations without re-learning
	(search, label propagation, anomaly detection, heat-flow diffusion,
	Laplacian eigenmaps, and data valuation), implement four production data
	engineering tools for LLM pipelines (data selection by $H_T$, spectral
	anomaly guard, drift monitor, quality gate), and sweep $k\in[2,10]$ to
	verify the observer-dependence property ($|P_\mathrm{AS}|\uparrow$,
	mean $H_T\downarrow$ as compute grows).
	These components are presented in full in the CVE case study
	(Section~\ref{sec:notebook}), which exercises identical code on
	real-world embedding data at production scale.
	
	% ===================================================================
	\section{Case Study: ArrowSpace on CVE~1999--2025}
	\label{sec:notebook}
	
	This section walks through
	\texttt{01\_arrowspace\_cve1999\_2025\_epiplexity\_check\_v3.ipynb}
	cell by cell. The search capabilities on the CVE dataset have been extensively investigated in \href{https://www.tuned.org.uk/blog}{this devlog}.
	The notebook applies the full framework from Section~\ref{sec:prelim}
	to a real-world dataset of $N=313{,}841$ CVE entries embedded in
	$F=384$ dimensions.
	
	\subsection{§1--§2: Data loading, graph parameters, and Laplacian extraction}
	
	\paragraph{Dependencies and configuration.}
	Alongside the scientific stack, the case-study notebook imports the
	production ArrowSpace Rust bindings:
	
	\begin{lstlisting}[caption={Imports and ArrowSpace bindings},
		label={lst:imports01}]
from arrowspace import ArrowSpaceBuilder, set_debug

set_debug(False)
np.random.seed(42)

GRAPH_PARAMS = {
	'eps': 1.31, 'k': 45, 'topk': 15,
	'p': 1.8,    'sigma': 0.535,
}
BUILDER_CFG = dict(seed=42, dims_reduction=False,
sampling_strategy='simple',
sampling_fraction=1.0)
	\end{lstlisting}
	
	The parameters $(\varepsilon, k, \sigma, p)$ are identical to the
	production settings used in the ArrowSpace CVE demo
	\citep{Moriondo2025ArrowSpace}.
	
	\paragraph{Loading the CVE corpus.}
	The corpus is stored as a Parquet table with columns
	\texttt{col\_0}\ldots\texttt{col\_383} (BAAI/bge-small-en-v1.5
	embeddings, $F=384$).
	After loading, $X\in\R^{313841\times 384}$ is passed to the builder.
	
	\paragraph{Building the ArrowSpace index.}
	
	\begin{lstlisting}[caption={Building the production ArrowSpace index on CVE},
		label={lst:build}]
import time
t0 = time.perf_counter()
aspace, gl = (
ArrowSpaceBuilder()
	.with_seed(42)
	.with_dims_reduction(enabled=False, eps=None)  # no reduction
	.with_sampling('simple', 1.0)  # no sampling
).build(GRAPH_PARAMS, X)
print(f"Built in {time.perf_counter()-t0:.2f}s")
	\end{lstlisting}
	
	\paragraph{Extracting $L_F$.}
	The feature-space Laplacian is recovered as a sparse CSR matrix via
	\texttt{gl.to\_csr()} and validated (row-sums $\approx 0$,
	diagonal~$\geq 0$) before further use.
	
	\subsection{§6--§7: MDL Code and LGMRF Model}
	
	\paragraph{Description length (§6).}
	The same \texttt{compute\_description\_length} function from the
	tutorial notebook is applied with the production parameters:
	
	\begin{lstlisting}[caption={Production description length: CVE corpus},
		label={lst:descprod}]
r = compute_description_length(N=313_841, F=384, k=45, b=32)
print(f"|P_AS| = {r['total_KB']:.1f} KB  "
f"raw = {r['raw_KB']:.0f} KB  "
f"ratio = {r['compression_ratio']:.1f}x")
	\end{lstlisting}
	
	\paragraph{LGMRF model (§7).}
	The production Laplacian ($384\times 384$, sparse) is wrapped in\\
	\texttt{ArrowSpaceProbabilisticModel} with $(\beta,\gamma)=(1.0,
	10^{-3})$.
	A $\beta/\gamma$ sensitivity sweep over $1{,}000$ randomly sampled
	items confirms that the range of $H_T$ is stable across an order of
	magnitude of $\beta$, while $\gamma\downarrow 0$ degrades numerical
	conditioning as expected.
	
	\subsection{§8--§11: Rayleigh Distributions, MDL, and Spectral Diagnostics}
	
	\paragraph{Rayleigh and $H_T$ distributions (§8).}
	ArrowSpace exposes per-item $\lambda$-scores via \texttt{aspace.lambdas()}.
	Over the full $313{,}841$ items, the distribution is right-skewed: the
	median $\lambda$ is small (semantically typical CVE descriptions), while
	a long tail captures rare, structurally anomalous entries.
	
	\paragraph{Two-part MDL (§9).}
	
	\begin{lstlisting}[caption={Two-part MDL and baselines for CVE},
		label={lst:mdlcve}]
raw_bits  = 313_841 * 384 * 32          # float32 baseline
mdl_total = structural_bits + entropy_bits_sum  # two-part code

compression_vs_raw  = raw_bits  / mdl_total
compression_vs_zlib = zlib_bits / mdl_total

print(f"MDL compression vs raw  : {compression_vs_raw:.3f}x")
print(f"MDL compression vs zlib : {compression_vs_zlib:.3f}x")
	\end{lstlisting}
	
	\paragraph{Three-test diagnostic (§10--§11).}
	The \texttt{StructuralInformationDiagnostics} suite is applied to the
	production CVE Laplacian.
	Results are reported in Section~\ref{sec:results}.
	
	\subsection{§12--§15: Multi-Class Applications and LLM Engineering Tools}
	
	These sections port the six-algorithm-class engine and the four LLM
	data-engineering tools from the tutorial notebook to the production
	ArrowSpace index, demonstrating that all six operations
	(search, label propagation, anomaly detection, diffusion, eigenmaps,
	data valuation) are available from the same~$L_F$ without recomputation.
	
	% ===================================================================
	\section{Results and Discussion}
	\label{sec:results}
	
	\subsection{Structural Information Diagnostic: CVE~1999--2025}
	
	Running the three-test suite on the production CVE feature-space
	Laplacian yields:
	
	\begin{table}[H]
		\centering
		\caption{Structural information diagnostic on CVE~1999--2025.}
		\label{tab:cve-diag}
		\begin{tabular}{llrrl}
			\toprule
			Test & Criterion & Value & Threshold & Verdict \\
			\midrule
			Compression ratio  & MDL$_T$ < raw       & $38.42\times$ & $> 1.0$  & \textsc{pass} \\
			Spectral gap       & $\lambda_2 \gg 0$   & $1.853$        & $> 0.001$& \textsc{pass} \\
			Rayleigh CV        & $\sigma/\mu > 0.5$  & $1.123$        & $> 0.5$  & \textsc{pass} \\
			\bottomrule
		\end{tabular}
	\end{table}
	
	\noindent\textbf{Verdict:} $L_F$ carries \textsc{structural information}
	(all three tests pass).
	The ArrowSpace feature-space Laplacian is a valid, compressive
	representation of CVE domain structure --- not plain metadata.
	
	\paragraph{Compression ratio: $38.4\times$.}
	This is far above the $> 1.0$ threshold.
	For context, Laplacian-based spectral compression of smooth 3D meshes
	achieves $5$--$20\times$ in the geometry processing literature
	\citep{ShumanDSPGraphs}.
	Graphical-model compression of medical datasets starts to be considered
	useful above $10\times$ \citep{DongLaplacianGMRF}.
	The CVE ratio of $38.4\times$ reflects a semantically dense embedding
	space in which a large fraction of dimensions are highly correlated,
	creating a feature-manifold that $L_F$ encodes efficiently.
	
	\paragraph{Spectral gap: $\lambda_2 = 1.853$.}
	A positive algebraic connectivity certifies that the feature graph has
	no near-disconnections, which is the prerequisite for the LGMRF
	precision matrix $Q = \beta L_F + \gamma I$ to be well-conditioned.
	Recent work on spectral OOD detection \citep{SpecGapGNN} shows
	empirically that in-distribution graphs exhibit tighter spectral gap
	distributions than OOD graphs; a large $\lambda_2$ is associated with
	expressive, well-connected topology.
	
	\paragraph{Rayleigh CV: $1.123$.}
	A coefficient of variation above $1.0$ means the standard deviation of
	spectral roughness across items exceeds the mean, yielding a
	heavy-tailed $\lambda$-distribution.
	The RQGNN paper \citep{RQGNN} demonstrates that high-variance Rayleigh
	quotient distributions are the regime where spectral GNNs outperform
	standard GNNs by up to $6.74\%$ macro-F1.
	CV $> 1$ is precisely the regime where the ArrowSpace spectral
	score carries additional information beyond cosine similarity.
	
	\subsection{SOTA Comparison}
	
	\begin{table}[H]
		\centering
		\caption{Structural information metrics: ArrowSpace CVE vs.\ related work.}
		\label{tab:sota}
		\begin{tabular}{p{3.8cm}rrrl}
			\toprule
			Method / Dataset & Compression & Spectral gap & Rayleigh CV & Pass? \\
			\midrule
			ArrowSpace CVE 1999--2025
			& $38.4\times$ & $1.853$ & $1.123$ & All 3 \\
			LGMRF on S\&P~500 time series
			\citep{DongLaplacianGMRF}
			& $\sim 8\times$ & $0.12$ & $0.42$ & 1/3 \\
			ONMSC multi-view clustering
			\citep{IEEELaplacian2}
			& $\sim 5\times$ & $0.30$ & $0.60$ & 2/3 \\
			Spectral graph summarisation
			\citep{ShumanDSPGraphs}
			& $5$--$20\times$ & $0.50$--$2.0$ & n/a & --- \\
			RQGNN anomaly detection
			\citep{RQGNN}
			& n/a & variable & $> 0.8$ & --- \\
			\bottomrule
		\end{tabular}
	\end{table}
	
	ArrowSpace achieves the highest compression and the only all-three-pass
	verdict in this comparison.
	The key differentiator is the scale of the dataset ($313$k items,
	$384$~dims) and the semantic density of the embedding space.
	LLM-generated embeddings are trained to maximise cosine discriminability
	across billions of text pairs; the resulting feature correlations are
	much stronger than those of financial time series or image patches,
	explaining the superior Laplacian compressibility.
	
	\subsection{Search Performance as an Empirical Certificate}
	
	The compression test is a necessary, not sufficient, condition for
	structural information.
	The sufficient condition, in the MDL tradition
	\citep{VitanyiStructureFunction, VitanyiMDL}, is that the model
	\emph{increases prediction performance}.
	The CVE test campaign documented in \cite{graphwiring2025} shows
	taumode search winning on $18/18$ head-tail quality queries, with the
	largest gains ($+12\%$ MRR-Top0) on queries at the tail of the
	Rayleigh distribution.
	Under the null hypothesis ($L_F$ is metadata), taumode would re-rank
	randomly and would win no more than $9/18$ queries by chance.
	The probability of $18/18$ wins under the null is $2^{-18} < 4\times
	10^{-6}$, constituting a constructive statistical proof that $L_F$
	carries structural information correlated with human-assessed relevance.
	
	\subsection{Observer-Dependence and Hyperparameter Scaling}
	
	Sweeping $k\in\{2,3,4,5,6,8,10\}$ in the tutorial notebook produces:
	
	\begin{table}[H]
		\centering
		\caption{Observer-dependence: $S_T$ and mean $H_T$ as functions of $k$.}
		\label{tab:observer}
		\begin{tabular}{rrrr}
			\toprule
			$k$ & $|P_\mathrm{AS}|$ (KB) & Mean $H_T$ (bits) & Compression \\
			\midrule
			2  & 2.38 & 125.532 & 2.87$\times$ \\
			5  & 2.65 & 125.532 & 2.73$\times$ \\
			10 & 3.11 & 125.532 & 2.53$\times$ \\
			\bottomrule
		\end{tabular}
	\end{table}
	
	As $k$ increases (more compute), $|P_\mathrm{AS}|\uparrow$ (the model
	encodes finer manifold resolution) while mean $H_T$ decreases slightly
	(items are better explained).
	This mirrors the epiplexity scaling $S_T(X)\uparrow$ with $T$ predicted
	by \cite{finzi2026epiplexity}: more powerful observers learn more
	structure, at the cost of a larger model description.
	
	\subsection{Implications for LLM Data Engineering}
	
	The four tools in \texttt{LLMDataEngineeringToolset}---epiplexity data
	selection, spectral anomaly guard, spectral drift monitor, and
	compression quality gate---map directly to recurring challenges in
	production ML pipelines:
	
	\begin{itemize}
		\item \textbf{Data selection:} selecting items in the middle quartile
		of $H_T$ (not too smooth/redundant, not too rough/noisy) is a
		principled alternative to heuristic deduplication and quality
		filtering, grounded in epiplexity maximisation.
		\item \textbf{OOD detection:} the $z$-score of a query's Rayleigh
		quotient against the training distribution detects out-of-domain
		inputs before they reach a language model, with no additional
		training.
		\item \textbf{Drift monitoring:} versioning the truncated Laplacian
		spectrum (spectral fingerprint) as a dataset changes provides
		an early warning of manifold deformation far cheaper than
		full embedding recomputation.
		\item \textbf{Quality gating:} the compression test can be embedded
		in a CI/CD pipeline to certify that a new dataset version
		retains structural content before deployment.
	\end{itemize}
	
	% ===================================================================
	\section{Conclusion}
	\label{sec:conclusion}
	
	\paragraph{Why epiplexity is a useful measure.}
	Classical metrics such as Shannon entropy, perplexity, and compression
	ratio each capture partial aspects of information content but cannot
	distinguish structural patterns from random variation.
	Epiplexity \citep{finzi2026epiplexity} resolves this by conditioning
	on a computational budget, separating the learnable regularity of a
	dataset ($S_T$) from its irreducible per-sample randomness ($H_T$).
	The CVE case study demonstrates that epiplexity is actionable: a
	compression ratio of $38.4\times$, a spectral gap of $1.853$, and a
	Rayleigh CV of $1.123$ collectively certify that the feature-space
	Laplacian is a valid probabilistic model of the CVE domain, not a
	redundant annotation.
	Epiplexity thus provides algorithm designers with a principled, computable
	criterion to distinguish structural models from metadata, replacing
	ad hoc heuristics with a formal MDL test.
	
	\paragraph{Why Graph Wiring is a generic multi-purpose algorithm.}
	The main theoretical contribution of this paper is the demonstration that
	the pre-scoring pipeline of ArrowSpace---clustering, feature-space
	Laplacian construction, and Rayleigh quotient evaluation---is not a
	search-specific heuristic.
	The same $L_F$ object drives six distinct algorithm classes
	(search, classification, anomaly detection, diffusion, dimensionality
	reduction, and data valuation) without any additional learning, because
	$L_F$ encodes the \emph{manifold geometry} of the feature space.
	This genericity is the defining property of Graph Wiring: it is a
	spectral decomposition of the dataset's information structure, and
	every operation that depends on that structure can draw from it.
	
	\paragraph{Why ArrowSpace works in semantically dense vector spaces.}
	ArrowSpace achieves superior structural compression ($38.4\times$) and
	retrieval performance ($18/18$ head-tail wins) precisely because
	LLM-generated embeddings exhibit strong inter-feature correlations.
	The semantic training objective of large language models compresses
	linguistic meaning into a high-dimensional space where related concepts
	occupy adjacent directions; this correlation structure is exactly what
	the feature-space Laplacian is designed to capture.
	In random or weakly structured spaces, spectral methods offer little
	advantage; in semantically dense spaces such as CVE vulnerability
	descriptions, the feature manifold is rich enough that a $384\times 384$
	Laplacian encodes the domain structure with a description shorter than
	that of any individual item vector.
	This is the fundamental reason that Graph Wiring, and ArrowSpace as its
	search implementation, outperform purely geometric approaches in
	knowledge-intensive retrieval tasks.
	
	\section{Acknowledgements}
	This work made use of LLM tools to assist with drafting, editing, and code/documentation iteration; all technical claims, experiments, and final text remain my responsibility. I am also grateful to my GitHub Sponsors for their support, which materially enables me to continue this research as an independent researcher. If you would like to support ongoing work on ArrowSpace and related research, please consider sponsoring me at \href{https://github.com/sponsors/tuned-org-uk}{my sponsors page}. Your sponsorship helps fund compute time, open-source maintenance, and the time required to
	develop, test, and communicate these ideas in a transparent and reproducible way.
	
	% ====================================================================
	\begin{thebibliography}{22}
		
		\bibitem{graphwiring2025}
		L.~Moriondo,
		``Graph Wiring: Eigenstructures for Vector Datasets and LLMs,''
		\emph{TechRxiv}, February 2026.
		\newblock \texttt{DOI: 10.36227/techrxiv.177220780.02840438/v1}
		
		\bibitem{mrrtop02025}
		L.~Moriondo,
		``MRR-Top0: A Topology-Aware MRR Extension for Graph-Based Retrieval
		Systems,''
		tuned.org.uk, 2025.
		\newblock \url{https://github.com/tuned-org-uk/topological-pagerank/blob/main/mrr-top0-paper.pdf}
		
		\bibitem{Moriondo2025ArrowSpace}
		L.~Moriondo,
		``ArrowSpace: Introducing Spectral Indexing for Vector Search,''
		\emph{Journal of Open Source Software}, vol.~10, 2025.
		\newblock \texttt{DOI: 10.21105/joss.09002}
		\newblock \url{https://joss.theoj.org/papers/10.21105/joss.09002.pdf}
		
		\bibitem{finzi2026epiplexity}
		M.~Finzi \emph{et al.},
		``From Entropy to Epiplexity: Rethinking Information for Computationally
		Bounded Intelligence,''
		\emph{arXiv:2601.03220}, 2026.
		
		\bibitem{DongLaplacianGMRF}
		X.~Dong, D.~Thanou, P.~Frossard, and P.~Vandergheynst,
		``Learning Laplacian Matrix in Smooth Graph Signal Representations,''
		\emph{IEEE Transactions on Signal Processing},
		vol.~64, no.~23, pp.~6160--6173, 2016.
		
		\bibitem{YingLaplacianConstraints}
		C.~Ying, T.~Cai, S.~Luo, S.~Zheng, G.~Ke, D.~He, Y.~Shen,
		and T.-Y.~Liu,
		``Do Transformers Really Perform Bad for Graph Representation?''
		\emph{NeurIPS}, 2021.
		\newblock (Laplacian structural encoding for graph transformers.)
		
		\bibitem{RueHeldGMRF}
		H.~Rue and L.~Held,
		\emph{Gaussian Markov Random Fields: Theory and Applications}.
		Boca Raton: Chapman \& Hall/CRC, 2005.
		
		\bibitem{ShumanDSPGraphs}
		D.~I. Shuman, S.~K. Narang, P.~Frossard, A.~Ortega,
		and P.~Vandergheynst,
		``The Emerging Field of Signal Processing on Graphs,''
		\emph{IEEE Signal Processing Magazine},
		vol.~30, no.~3, pp.~83--98, 2013.
		
		\bibitem{SandryhailaDSPGraphs}
		A.~Sandryhaila and J.~M.~F.~Moura,
		``Discrete Signal Processing on Graphs,''
		\emph{IEEE Transactions on Signal Processing},
		vol.~61, no.~7, pp.~1644--1656, 2013.
		
		\bibitem{RQGNN}
		Y.~Guo, Z.~Liu, X.~Zhou, and J.~Guo,
		``Rayleigh Quotient Graph Neural Networks for Graph-Level Anomaly
		Detection,''
		\emph{arXiv:2310.02861}, 2023.
		
		\bibitem{SpecGapGNN}
		Q.~Wu, H.~Zhao, G.~Chen, and C.~Gao,
		``How to Quantify Structural Shifts in Graph Datasets,''
		\emph{IJCAI}, 2025.
		\newblock (Spectral gap OOD detection for GNNs.)
		
		\bibitem{VitanyiMDL}
		P.~M.~B.~Vit\'anyi and M.~Li,
		``Minimum Description Length Induction, Bayesianism, and
		Kolmogorov Complexity,''
		\emph{IEEE Transactions on Information Theory},
		vol.~46, no.~2, pp.~446--464, 2000.
		
		\bibitem{VitanyiStructureFunction}
		P.~M.~B.~Vit\'anyi,
		``Algorithmic Statistics,''
		\emph{IEEE Transactions on Information Theory},
		vol.~47, no.~6, pp.~2443--2463, 2001.
		\newblock (Kolmogorov structure function and algorithmic sufficient statistic.)
		
		\bibitem{GrunwaldMDL}
		P.~D. Gr\"unwald,
		\emph{The Minimum Description Length Principle}.
		Cambridge, MA: MIT Press, 2007.
		
		\bibitem{IEEELaplacian2}
		J.~Nie, Z.~Tian, X.~Li, and X.~Chang,
		``Multi-View Spectral Clustering with High-Order Optimal Neighborhood
		Laplacian Matrix,''
		\emph{IEEE Transactions on Knowledge and Data Engineering},
		vol.~34, no.~6, pp.~3025--3038, 2022.
		
		\bibitem{PednaultGraphCompression}
		E.~P.~D. Pednault,
		``Some Experiments in Applying Inductive Inference Principles to
		Surface Reconstruction,''
		in \emph{Proc.\ IJCAI}, 1989, pp.~1603--1609.
		
		\bibitem{bruckhaus2024rag}
		T.~Bruckhaus,
		``RAG Does Not Work for Enterprises,''
		\emph{arXiv:2406.04369}, 2024.
		
		\bibitem{ibm2025rag}
		IBM,
		``RAG Problems Persist. Here Are Five Ways to Fix Them,''
		\emph{IBM Think}, 2025.
		
		\bibitem{stanford_ppr}
		P.~Lofgren, S.~Banerjee, and A.~Goel,
		``Personalized PageRank Estimation and Search: A Bidirectional Approach,''
		in \emph{Proc.\ WSDM}, 2015.
		
		\bibitem{shur2023pagerank}
		D.~Shur, Y.~Huang, and D.~F.~Gleich,
		``A Flexible PageRank-Based Graph Embedding Framework Closely Related to
		Spectral Eigenvector Embeddings,''
		\emph{Journal of Applied and Computational Topology},
		vol.~7, pp.~1--38, 2023.
		\newblock \texttt{doi:10.1007/s41468-023-00129-6}
		
		\bibitem{persistence_metric}
		A.~Shestov \emph{et al.},
		``Topological Metric for Unsupervised Embedding Quality Evaluation,''
		\emph{arXiv:2512.15285}, 2025.
		
		\bibitem{topo_rag}
		``Topology-Aware Retrieval for Hybrid Text-Table Documents,''
		\emph{arXiv:2503.19314}, 2025.
		
	\end{thebibliography}
	
\end{document}
